\documentclass{article}
\usepackage[margin=1in]{geometry}
\usepackage{mathtools, amsfonts, amsthm, xcolor, listings}

\title{HW04}
\author{Sam Ly}


\begin{document}
\maketitle

\section*{Exercise 1 [2]}
\noindent
Tools and resources used: Mainly Google, and Gemini when I got really stuck.
\begin{enumerate}
  \item[Exercise 1]
    {\color{red} Points attempted: 2}.
    {\color{blue} Time spent: 10 minutes}.
    {\color{brown} Difficulty: 1/10}.
    {\color{green!50!black} Self-assessed score: 2/2}.
    {\color{magenta} Questions/Feedback: None.}
  \item[Exercise 2]
    {\color{red} Points attempted: 7}.
    {\color{blue} Time spent: 18 minutes}.
    {\color{brown} Difficulty: 7/10}.
    {\color{green!50!black} Self-assessed score: 3/3}.
    {\color{magenta} Questions/Feedback: How should I be thinking about gradients?}
\end{enumerate}
\pagebreak

\section*{Exercise 2 [7]}
\begin{enumerate}
    \item {
        We defined the dihedral group in terms of its group presentation: \[         D_n = \langle r, f \mid r^n = f^2 = (rf)^2 = e \rangle.       \]
      
        \begin{enumerate}
            \item {
                Prove that \(r^kf = fr^{n-k}\) for all \(0 < k < n\).

                \begin{proof}
                    First, we notice that \(rfrf = e\). By right muliplying by 
                    \(f\) we get \(rfr = f\). By left muliplying by \(r^{-1}\), 
                    we get \(fr = r^{-1}f\). Then, we left and right muliply by 
                    \(f\) on both sides to get \(rf = fr^{-1}f\).

                    In other words, \(r^1f = fr^{n-1}\). Now, to prove that 
                    \(r^kf = fr^{n-k}\) for all \(0 < k < n\), we must use 
                    induction. We have already proved the base case \(k = 1\). 
                    As our inductive hypothesis, we assume that the proposition 
                    is true for some value \(k\). Thus, \(r^kf = fr^{n-k}\).
                    Now, for by left multiplying by \(r\), we get 
                    \[r^{k+1}f = rfr^{n-k}.\]

                    Now, we apply \(rf = fr^{n-1}\) to the RHS, and get 
                    \[r^{k+1}f = fr^{-1}r^{n-k} = fr^{n- (k+1)}.\]
                    
                    Therefore, \(r^kf = fr^{n-k}\) for all \(0 < k < n\).
                \end{proof}
            }
            \item {
                Prove that every element can be written in the form \(r^kf^\ell\) 
                for \(0 \leq k < n\) and \(\ell \in \{0, 1\}\)

                \begin{proof}
                    We approach this problem from  a formal languages angle. 
                    We begin by constructing all possible `words' of the language 
                    formed by the alphabet \(\{r, s \}\). One example of such a 
                    word is \(rrrfrfffrffrrrrrrrf\). By definition, this language
                    must encompass all of \(D_n\) for any \(n\). 
                    
                    Then, we can form some 
                    basic substitution rules, in other words, a `grammar' that allows 
                    us to rewrite the word into a nicer form. By definition of the 
                    dihedral group, we have the rules: 
                    \begin{equation}
                        \underbrace{rr \dots r}_{m \text{ terms}} = r^m = r^{m-n}
                    \end{equation}
                    \begin{equation}
                        ff = e. \\
                    \end{equation}

                    Thus, we can rewrite our word into the form 
                    \(r^{k_0}f^{l_0} r^{k_1} f^{l_1} \dots \),
                    where \(0 \le k_i < n\) and \(l_i \in \{0, 1\}\). 

                    Now, from here, we apply the proposition proved above to 
                    combine and move all \(l_i\) to the right of the expression.

                    For example, if \(l_0 = 0\), then the expression becomes 
                    \(r^{k_0}r^{k_1} f^{l_1} \dots = r^{k_0 + k_1} f^{l_1} \dots \).
                    If \(l_0 = 1\), then the expression becomes 
                    \[r^{k_0}(fr^{k_1})f^{l_1} \dots = r^{k_0} (r^{n - k_1}f) f^{l_1} \dots \]
                    which we can simplify by combining \(ff^{l_1}\). 
                    This process can be repeated until we reach the form \(r^kf^\ell\) 
                    for \(0 \leq k < n\) and \(\ell \in \{0, 1\}\).
                \end{proof}
            }
            \item {
                Prove that the order (number of elements) of the group is \(D_n = 2n\).

                \begin{proof}
                    Now that we have established that all elements of \(D_n\) can 
                    be written as the form \(r^kf^\ell\) 
                    for \(0 \leq k < n\) and \(\ell \in \{0, 1\}\), this queestion
                    can be reframed as asking ``How many unique ways can we write 
                    \(r^kf^\ell\)?'' By combinatorics, we see that the number of 
                    unique ways is \(2n\).
                \end{proof}
            }

        \end{enumerate}
    }

    \item {
        We say that a group \((G,*)\) is an abelian group if and only if 
        \(a * b = b * a\) for all \(a, b \in G\).
      
        \begin{enumerate}
            \item {
                Prove that if \(G\) is generated by \(g_1\) and \(g_2\) then 
                \(G\) is abelian if and only if \(g_1 * g_2 = g_2 * g_1\).


                \begin{proof}
                    First, we prove that if \(a * b = b * a\) for all \(a, b \in G\),
                    then \(g_1 * g_2 = g_2 * g_1\). By setting \(a = g_1\) and 
                    \(b = g_2\),  then \(g_1 * g_2 = g_2 * g_1\).

                    The converse ``if \(g_1 * g_2 = g_2 * g_1\), then \(a * b = b * a\) for all \(a, b \in G\)''
                    is much more difficult to prove. 

                    % First, we prove that \({g_1}^kg_2 = g_2g^k\). We do this by seeing 
                    % \[{g_1}^kg_2 = {g_1}^{k-1} g_1 g_2 = {g_1}^{k-1}g_2g_1.\]

                    % Then, notice that \({g_1}^{k-1}g_2 = {g_1}^{k-2} g_1 g_2\), thus 
                    % \[{g_1}^{k-1}g_2g_1 = {g_1}^{k-2} g_2 {g_1}^2.\] 

                    % In general, for some \(x \le k\), 
                    % \[{g_1}^kg_2 = {g_1}^{k-x}g_2{g_1}^x.\]

                    % Now, we can set \(x = k\) times to see that 
                    % \({g_1}^kg_2 = g_2{g_1}^k\).
                    
                    % With this lemma proven, we can show that 
                    % \({g_1}^k {g_2}^l = {g_2}^l {g_1}^k\).

                    % First, we see that 
                    % \begin{align}
                    % {g_1}^k {g_2}^l &= {g_1}^{k-1} (g_1 {g_2}^l) \\
                    %                 &= {g_1}^{k-1} ({g_2}^l g_1) \\
                    %                 &= {g_1}^{k-x} {g_2}^l {g_1}^x \\
                    %                 &= {g_2}^l {g_1}^k .
                    % \end{align}

                    % Now, all of the machinery is in place to prove that 
                    % \(a * b = b * a\) for all \(a, b \in G\).

                    Notice that 
                    \begin{align}
                        a &= a_1 * a_2 * \dots * a_p = \prod_{i=1}^p a_i \\
                        b &= b_1 * b_2 * \dots * b_q = \prod_{i=1}^q b_i \\
                        c &= a * b = a_1 * a_2 * \dots * a_p * b_1 * b_2 * \dots * b_q
                    \end{align}
                    where \(a_i, b_i \in \{g_1, g_2\}\).

                    Notice that because \(g_1 * g_2 = g_2 * g_1\), we can swap 
                    any two adjacent terms of \(c\).

                    Therefore, we can rewrite \(c\) as 
                    \begin{align}
                        c &= a * b = b_1 * b_2 * \dots * b_q * a_1 * a_2 * \dots * a_p.
                    \end{align}

                    Therefore, \(G\) is generated by \(g_1\) and \(g_2\) then 
                    \(G\) is abelian if and only if \(g_1 * g_2 = g_2 * g_1\).
                \end{proof}
            }
        
        
        \end{enumerate}
    }
\end{enumerate}

\end{document}
